
\documentclass{article}
\usepackage{colortbl}
\usepackage{makecell}
\usepackage{multirow}
\usepackage{supertabular}

\begin{document}

\newcounter{utterance}

\centering \large Interaction Transcript for game `cladder', experiment `full\_v1.5\_default', episode 9786 with Qwen3.6{-}27B.
\vspace{24pt}

{ \footnotesize  \setcounter{utterance}{1}
\setlength{\tabcolsep}{0pt}
\begin{supertabular}{c@{$\;$}|p{.15\linewidth}@{}p{.15\linewidth}p{.15\linewidth}p{.15\linewidth}p{.15\linewidth}p{.15\linewidth}}
    \# & \multicolumn{2}{c}{Player} && \multicolumn{2}{c}{Game Master} \\
    \hline

    \theutterance \stepcounter{utterance}  
    & & & \multicolumn{4}{p{0.6\linewidth}}{
        \cellcolor[rgb]{0.9,0.9,0.9}{
            \makecell[{{p{\linewidth}}}]{
                \texttt{\tiny{[P1$\langle$GM]}}
                \texttt{You are an expert in causal inference. The following question is not a typical commonsense query, but rather a meticulously designed question created by a professor specializing in causal inference, intended to assess the students' mastery of the course content.} \\
\texttt{Imagine a self{-}contained, hypothetical world with only the following conditions, and without any unmentioned factors or causal relationships: Qwiu has a direct effect on jyka and kwox. Jyka has a direct effect on yupt. Yupt has a direct effect on kwox. Qwiu is unobserved. We know that qwiu causes jyka. jyka causes yupt. qwiu or yupt causes kwox. We observed an individual is not qwiu. Would an individual is not kwox if not jyka instead of jyka?} \\
\texttt{Start your answer with 'Yes' or 'No', followed by additional reasoning or evidence to support your explanation.} \\
            }
        }
    }
     \\ \\

    \theutterance \stepcounter{utterance}  
    & \multicolumn{4}{p{0.6\linewidth}}{
        \cellcolor[rgb]{0.561,0.737,0.561}{
            \makecell[{{p{\linewidth}}}]{
                \texttt{\tiny{[P1$\rangle$GM]}}
                \texttt{No} \\
\\ 
\texttt{**Reasoning:**} \\
\\ 
\texttt{To determine the answer, we must reconstruct the causal structure and evaluate the counterfactual query based on the provided conditions.} \\
\\ 
\texttt{1.  **Causal Structure Identification:**} \\
\texttt{    *   \$Q \textbackslash{}to J\$ (Qwiu causes Jyka)} \\
\texttt{    *   \$J \textbackslash{}to Y\$ (Jyka causes Yupt)} \\
\texttt{    *   \$Q \textbackslash{}to K\$ (Qwiu causes Kwox)} \\
\texttt{    *   \$Y \textbackslash{}to K\$ (Yupt causes Kwox)} \\
\texttt{    *   The prompt states "Qwiu or Yupt causes Kwox," which in causal inference contexts with direct effects listed usually implies that both are direct causes of Kwox (i.e., \$K\$ is a function of \$Q\$ and \$Y\$). The structure is a "fork" at \$Q\$ and a "chain" \$Q \textbackslash{}to J \textbackslash{}to Y\$, with \$Y\$ and \$Q\$ both pointing to \$K\$.} \\
\texttt{    *   \$Q\$ is unobserved.} \\
\texttt{    *   We observe an individual where \$Q=0\$ (not Qwiu).} \\
\\ 
\texttt{2.  **The Counterfactual Query:**} \\
\texttt{    *   The question asks: "Would an individual be not Kwox if not Jyka instead of Jyka?"} \\
\texttt{    *   This is a counterfactual query: \$K\_\{J=0\} =} \\
            }
        }
    }
    & & \\ \\

    \theutterance \stepcounter{utterance}  
    & & & \multicolumn{2}{p{0.3\linewidth}}{
        \cellcolor[rgb]{0.9,0.9,0.9}{
            \makecell[{{p{\linewidth}}}]{
                \texttt{\tiny{[GM$|$GM]}}
                \texttt{No} \\
            }
        }
    }
    & & \\ \\

    \theutterance \stepcounter{utterance}  
    & & & \multicolumn{2}{p{0.3\linewidth}}{
        \cellcolor[rgb]{0.9,0.9,0.9}{
            \makecell[{{p{\linewidth}}}]{
                \texttt{\tiny{[GM$|$GM]}}
                \texttt{yes} \\
            }
        }
    }
    & & \\ \\

    \theutterance \stepcounter{utterance}  
    & & & \multicolumn{2}{p{0.3\linewidth}}{
        \cellcolor[rgb]{0.9,0.9,0.9}{
            \makecell[{{p{\linewidth}}}]{
                \texttt{\tiny{[GM$|$GM]}}
                \texttt{game\_result = LOSE} \\
            }
        }
    }
    & & \\ \\

\end{supertabular}
}

\end{document}
